% GENERATED FILE. Edit canonical lessons, then run npm run generate:books.
% canonical-source-hash: fnv1a64:e41913104060e074
% canonical-lessons: GE-C352-entlassen, GE-C352-arbeitslos, GE-C352-arbeitslosigkeit, GE-C352-probezeit, GE-C352-einarbeiten

\chapter{Ohne Arbeit}
\label{ch:ge-ohne-arbeit}
\hlchaptermodality{352}

\begin{chapteropening}
\textbf{By the end of this chapter:} \emph{I can talk about losing a job and being out of work, and about a trial period and settling into a new job.}

Complete the last lesson of chapter 352: I can talk about losing a job and being out of work, and about a trial period and settling into a new job.
\end{chapteropening}

\section[entlassen]{entlassen — to lay off, to dismiss; to discharge}

\label{lesson:GE-C352-entlassen}



\begin{quote}
Before the new one: say the German for the hour of overtime, then the German for the end of the working day.
\end{quote}

\subsection*{You'll want to know: entlassen}
\textbf{entlassen} — \textquotedblleft{}to lay off, to dismiss; to discharge\textquotedblright{}.

\begin{quote}
Die Fabrik hat fünfzig Mitarbeiter entlassen.
\end{quote}

\begin{quote}
\emph{The factory laid off fifty employees.}
\end{quote}

Inseparable and strong: \emph{sie entlässt, sie hat entlassen}. From hospital: \emph{Er wird morgen entlassen} — he is being discharged tomorrow.

\begin{cousinweb}
\emph{entlassen} is built with \emph{ent-} \textquotedblleft{}away\textquotedblright{} on \emph{lassen} \textquotedblleft{}to let, to let go\textquotedblright{}, inherited from Old High German \emph{lāzan} and Proto-Germanic \emph{lētaną} (EtymDB): you let someone go.
\end{cousinweb}

\begin{grammarlens}[title={What you've built}]
One more word for telling, explaining and describing at B1.
\end{grammarlens}

\subsection*{Your turn}
\begin{itemize}
\raggedright
  \item \emph{Say it:} \emph{entlassen}
  \item \emph{Say it:} \emph{entlassen}, once more
  \item \emph{Say it:} a sentence of your own with \emph{entlassen}
  \item \emph{Recall:} say the German for the hour of overtime, then the German for the end of the working day, then say \emph{entlassen} again
\end{itemize}

\begin{tcolorbox}[breakable,colback=teal!4,colframe=teal!35!black,title={Before you move on}]
What does entlassen mean? (\textquotedblleft{}To lay off, to dismiss; to discharge\textquotedblright{}.) Say it once more.
\end{tcolorbox}
\section[arbeitslos]{arbeitslos — unemployed, out of work}

\label{lesson:GE-C352-arbeitslos}



\begin{quote}
Before the new one: say the German for the end of the working day, then the German for to lay off.
\end{quote}

\subsection*{You'll want to know: arbeitslos}
\textbf{arbeitslos} — \textquotedblleft{}unemployed, out of work\textquotedblright{}.

\begin{quote}
Nach der Kündigung war er drei Monate arbeitslos.
\end{quote}

\begin{quote}
\emph{After he was given notice, he was out of work for three months.}
\end{quote}

\emph{Der Arbeitslose, die Arbeitslose} — the unemployed person, declined like an adjective.

\begin{cousinweb}
\emph{arbeitslos} joins \emph{Arbeit} \textquotedblleft{}work\textquotedblright{}, Proto-Germanic \emph{arbaidiz} (EtymDB), to \emph{-los} \textquotedblleft{}without\textquotedblright{}, from Old High German \emph{lōs} \textquotedblleft{}loose, free, lacking\textquotedblright{} (EtymDB).
\end{cousinweb}

\begin{grammarlens}[title={What you've built}]
One more word for telling, explaining and describing at B1.
\end{grammarlens}

\subsection*{Your turn}
\begin{itemize}
\raggedright
  \item \emph{Say it:} \emph{arbeitslos}
  \item \emph{Say it:} \emph{arbeitslos}, once more
  \item \emph{Say it:} a sentence of your own with \emph{arbeitslos}
  \item \emph{Recall:} say the German for the end of the working day, then the German for to lay off, then say \emph{arbeitslos} again
\end{itemize}

\begin{tcolorbox}[breakable,colback=teal!4,colframe=teal!35!black,title={Before you move on}]
What does arbeitslos mean? (\textquotedblleft{}Unemployed, out of work\textquotedblright{}.) Say it once more.
\end{tcolorbox}
\section[die Arbeitslosigkeit]{die Arbeitslosigkeit — unemployment}

\label{lesson:GE-C352-arbeitslosigkeit}



\begin{quote}
Before the new one: say the German for to lay off, then the German for unemployed.
\end{quote}

\subsection*{You'll want to know: die Arbeitslosigkeit}
\textbf{die Arbeitslosigkeit} — \textquotedblleft{}unemployment\textquotedblright{}.

\begin{quote}
In dieser Region ist die Arbeitslosigkeit sehr hoch.
\end{quote}

\begin{quote}
\emph{Unemployment is very high in this region.}
\end{quote}

No plural in everyday use. \emph{Die Arbeitslosenquote} — the unemployment rate.

\begin{cousinweb}
\emph{Arbeitslosigkeit} is formed from \emph{arbeitslos} with \emph{-igkeit}; the adjective joins \emph{Arbeit} \textquotedblleft{}work\textquotedblright{} to \emph{-los} \textquotedblleft{}without\textquotedblright{}, from Old High German \emph{lōs} (EtymDB).
\end{cousinweb}

\begin{grammarlens}[title={What you've built}]
One more word for telling, explaining and describing at B1.
\end{grammarlens}

\subsection*{Your turn}
\begin{itemize}
\raggedright
  \item \emph{Say it:} \emph{die Arbeitslosigkeit}
  \item \emph{Say it:} \emph{die Arbeitslosigkeit}, once more
  \item \emph{Say it:} a sentence of your own with \emph{die Arbeitslosigkeit}
  \item \emph{Recall:} say the German for to lay off, then the German for unemployed, then say \emph{die Arbeitslosigkeit} again
\end{itemize}

\begin{tcolorbox}[breakable,colback=teal!4,colframe=teal!35!black,title={Before you move on}]
What does die Arbeitslosigkeit mean? (\textquotedblleft{}Unemployment\textquotedblright{}.) Say it once more.
\end{tcolorbox}
\section[die Probezeit]{die Probezeit — the trial period, the probation}

\label{lesson:GE-C352-probezeit}



\begin{quote}
Before the new one: say the German for unemployed, then the German for unemployment.
\end{quote}

\subsection*{You'll want to know: die Probezeit}
\textbf{die Probezeit} — \textquotedblleft{}the trial period, the probation\textquotedblright{}.

\begin{quote}
In der Probezeit kann man schneller kündigen.
\end{quote}

\begin{quote}
\emph{During the trial period you can give notice more quickly.}
\end{quote}

Plural: \emph{die Probezeiten}. Usually the first six months of a new job.

\begin{cousinweb}
\emph{Probezeit} joins \emph{Probe} \textquotedblleft{}trial, test\textquotedblright{}, borrowed from Late Latin \emph{proba} \textquotedblleft{}proof\textquotedblright{}, from \emph{probāre} \textquotedblleft{}to test, to approve\textquotedblright{} (EtymDB; Lewis \& Short), to \emph{Zeit} \textquotedblleft{}time\textquotedblright{}.
\end{cousinweb}

\begin{grammarlens}[title={What you've built}]
One more word for telling, explaining and describing at B1.
\end{grammarlens}

\subsection*{Your turn}
\begin{itemize}
\raggedright
  \item \emph{Say it:} \emph{die Probezeit}
  \item \emph{Say it:} \emph{die Probezeit}, once more
  \item \emph{Say it:} a sentence of your own with \emph{die Probezeit}
  \item \emph{Recall:} say the German for unemployed, then the German for unemployment, then say \emph{die Probezeit} again
\end{itemize}

\begin{tcolorbox}[breakable,colback=teal!4,colframe=teal!35!black,title={Before you move on}]
What does die Probezeit mean? (\textquotedblleft{}The trial period, the probation\textquotedblright{}.) Say it once more.
\end{tcolorbox}
\section[einarbeiten]{einarbeiten — to train in (a new colleague); to settle in (sich einarbeiten)}

\label{lesson:GE-C352-einarbeiten}



\begin{quote}
Before the new one: say the German for unemployment, then the German for the trial period.
\end{quote}

\subsection*{You'll want to know: einarbeiten}
\textbf{einarbeiten} — \textquotedblleft{}to train in (a new colleague); to settle in (sich einarbeiten)\textquotedblright{}.

\begin{quote}
Eine Kollegin arbeitet mich in den ersten Wochen ein.
\end{quote}

\begin{quote}
\emph{A colleague is training me in during the first weeks.}
\end{quote}

Separable: \emph{sie arbeitet mich ein}, \emph{sie hat mich eingearbeitet}. \emph{Sich einarbeiten} — to get to know a new job.

\begin{cousinweb}
\emph{einarbeiten} is built from \emph{ein} \textquotedblleft{}in\textquotedblright{} and \emph{arbeiten} \textquotedblleft{}to work\textquotedblright{}, from \emph{Arbeit}, Old High German \emph{arbeit} and Proto-Germanic \emph{arbaidiz} \textquotedblleft{}labour, hardship\textquotedblright{} (EtymDB): you work someone into a job.
\end{cousinweb}

\begin{grammarlens}[title={What you've built}]
One more word for telling, explaining and describing at B1.
\end{grammarlens}

\subsection*{Your turn}
\begin{itemize}
\raggedright
  \item \emph{Say it:} \emph{einarbeiten}
  \item \emph{Say it:} \emph{einarbeiten}, once more
  \item \emph{Say it:} a sentence of your own with \emph{einarbeiten}
  \item \emph{Recall:} say the German for unemployment, then the German for the trial period, then say \emph{einarbeiten} again
\end{itemize}

\begin{tcolorbox}[breakable,colback=teal!4,colframe=teal!35!black,title={Before you move on}]
What does einarbeiten mean? (\textquotedblleft{}To train in (a new colleague); to settle in (sich einarbeiten)\textquotedblright{}.) Say it once more.
\end{tcolorbox}
